\paragraph{玫瑰轨迹的 Lee--Yang 分离平方律}
沿用本小节的坐标门 \(J\)（\eqref{eq:lee-yang-J-def}）与 Joukowsky 映射 \(W(z)=z+z^{-1}\)，并沿用命题 \ref{prop:jouwkowsky-inverse-square} 的恒等式 \(J(z^2)=-1/W(z)^2\)。

\begin{theorem}[Lee--Yang 分离平方律]\label{thm:leyang-rose-separation-square-law}
令 \(n,d\in\NN\)。
对任意 \(w\in\CC^{\times}\) 且满足 \(w^{2n}\neq -1\)（从而 \(1+w^{2n}\neq 0\)），定义
\[
\Gamma_{n,d}(w):=\frac12\bigl(w^{d+n}+w^{d-n}\bigr),\qquad
\tau_n(w):=J(w^{2n})=-\frac{w^{2n}}{(1+w^{2n})^2}.
\]
则成立严格恒等式
\begin{equation}\label{eq:leyang-rose-separation-square-law}
\boxed{\ 4\,\tau_n(w)\,\Gamma_{n,d}(w)^2=-\,w^{2d}\ }.
\end{equation}
特别地，当 \(w\in\TT\) 时有 \(\tau_n(w)\in(-\infty,-\tfrac14]\)，并且
\[
\abs{\Gamma_{n,d}(w)}=\frac{1}{2\sqrt{-\tau_n(w)}},
\qquad
\arg\bigl(\Gamma_{n,d}(w)^2\bigr)\equiv 2d\,\arg(w)\pmod{2\pi}.
\]
\end{theorem}
\begin{proof}
由定义可写
\[
\Gamma_{n,d}(w)=\frac12\,w^{d-n}\,(1+w^{2n}),
\]
从而
\[
\Gamma_{n,d}(w)^2=\frac14\,w^{2d-2n}(1+w^{2n})^2.
\]
再与 \(\tau_n(w)=-\frac{w^{2n}}{(1+w^{2n})^2}\) 相乘即得
\(
\tau_n(w)\Gamma_{n,d}(w)^2=-\frac14 w^{2d}
\)，从而得到 \eqref{eq:leyang-rose-separation-square-law}。
当 \(w\in\TT\) 时，由命题 \ref{prop:jouwkowsky-inverse-square}（取 \(z=w^n\)）得 \(\tau_n(w)=J((w^n)^2)\in(-\infty,-\tfrac14]\)；
对 \eqref{eq:leyang-rose-separation-square-law} 取模与辐角即可得到所示两式。证毕。
\end{proof}

\begin{corollary}[关于 \(n\) 的守恒量]\label{cor:leyang-rose-n-invariant}
在定理 \ref{thm:leyang-rose-separation-square-law} 的记号下，令
\[
Q_{n,d}(w):=\tau_n(w)\,\Gamma_{n,d}(w)^2.
\]
则对一切 \(n\in\NN\) 恒有
\[
\boxed{\ Q_{n,d}(w)\equiv-\frac14\,w^{2d}\ }.
\]
因此，当 \(w\) 沿单位圆 \(\TT\) 绕行一周时，曲线 \(w\mapsto Q_{n,d}(w)\) 绕原点的净绕数严格等于 \(2d\)，与参数 \(n\) 无关。
\end{corollary}
\begin{proof}
由 \eqref{eq:leyang-rose-separation-square-law} 直接得到 \(Q_{n,d}(w)=-\frac14 w^{2d}\)。
当 \(w=e^{\mathrm{i}\theta}\)（\(\theta\in[0,2\pi]\)）时，\(w^{2d}\) 的辐角增量为 \(4\pi d\)，故绕数为 \(2d\)。证毕。
\end{proof}

\paragraph{Lissajous 的扭曲纤维积与消元闭式}
对 Lissajous 族 \eqref{eq:leyang-lissajous-family}，命题 \ref{prop:leyang-lissajous-unit-circle-uplift} 已给出二维 Lee--Yang 坐标 \(\Lambda_{a,b,\delta}\)。
本段改用“扭曲粘合”的代数语言刻画其像的 Zariski 闭包，并将该对象与正文的算术奇异环纤维积 \(\mathfrak{R}_S=\Sigma_S\times_{\TT}\TT^S\)（\S\ref{subsubsec:cdim-solenoid-arithmetic-singular-ring}）在结构上对齐为一种纯 Archimedean 两通道特例。

\begin{theorem}[扭曲纤维积的 Lee--Yang 消元方程]\label{thm:leyang-lissajous-twisted-fiberproduct-resultant}
令 \(a,b\in\ZZ_{>0}\) 且 \(\delta\in\RR\)，并记
\[
\xi:=e^{\mathrm{i}2b\delta}\in\TT.
\]
定义相位约束子集
\begin{equation}\label{eq:leyang-lissajous-Kabdelta}
K_{a,b,\delta}:=\{(u,v)\in\TT^2:\ u^b=\xi\,v^a\}.
\end{equation}
并定义其 Lee--Yang 像（扭曲纤维积）
\[
\mathcal{R}^{\mathrm{LY}}_{a,b,\delta}
:=\Bigl\{(T_1,T_2)\in\RR^2:\ \exists(u,v)\in K_{a,b,\delta},\ T_1=J(u),\ T_2=J(v)\Bigr\}.
\]
则 \(\mathcal{R}^{\mathrm{LY}}_{a,b,\delta}\subset(-\infty,-\tfrac14]^2\)，并且其 Zariski 闭包包含在代数曲线
\begin{equation}\label{eq:leyang-lissajous-resultant-curve}
F_{a,b,\delta}(T_1,T_2)=0
\end{equation}
之内，其中 \(F_{a,b,\delta}\in\CC[T_1,T_2]\) 可取为迭代结果式
\begin{equation}\label{eq:leyang-lissajous-resultant-def}
\boxed{\;
F_{a,b,\delta}(T_1,T_2)
:=\Res_v\!\Bigl(\Res_u\bigl(T_1(1+u)^2+u,\ u^b-\xi v^a\bigr),\ T_2(1+v)^2+v\Bigr)\; } .
\end{equation}
此外，\(F_{a,b,\delta}\) 对 \((T_1,T_2)\) 的双次数满足
\begin{equation}\label{eq:leyang-lissajous-resultant-bidegree-bound}
\boxed{\ \deg_{T_1}F_{a,b,\delta}\le 2b,\qquad \deg_{T_2}F_{a,b,\delta}\le 2a\ }.
\end{equation}
\end{theorem}
\begin{proof}
由 \(T=J(u)=-u/(1+u)^2\) 的等价变形得到二次门方程
\begin{equation}\label{eq:leyang-gate-quadratic}
T(1+u)^2+u=0,
\end{equation}
同理 \(T_2(1+v)^2+v=0\)。
若 \((T_1,T_2)\in\mathcal{R}^{\mathrm{LY}}_{a,b,\delta}\)，则存在 \((u,v)\in K_{a,b,\delta}\) 同时满足 \eqref{eq:leyang-gate-quadratic}（对 \(T_1,u\)）、其 \(v\)-版本以及 \(u^b-\xi v^a=0\)。
因此 \((T_1,T_2)\) 必使迭代结果式 \eqref{eq:leyang-lissajous-resultant-def} 为零，从而 \(\overline{\mathcal{R}^{\mathrm{LY}}_{a,b,\delta}}^{\,\mathrm{Zar}}\subseteq V(F_{a,b,\delta})\)。
双次数上界 \eqref{eq:leyang-lissajous-resultant-bidegree-bound} 由结果式的次数界给出：固定 \(v\) 时，\(\Res_u(\cdot)\) 对 \(T_1\) 的次数不超过 \(2b\)（二次门方程与 \(b\)-次扭曲幂约束消元），再与 \(v\)-门方程作一次结果式消元不增加 \(T_1\) 的次数；对 \(T_2\) 同理得到不超过 \(2a\)。证毕。
\end{proof}

\begin{proposition}[扭曲约束的连通分支数]\label{prop:leyang-lissajous-K-components}
设 \(g:=\gcd(a,b)\)，并写 \(a=ga_0\)、\(b=gb_0\) 且 \(\gcd(a_0,b_0)=1\)。
则集合 \(K_{a,b,\delta}\)（\eqref{eq:leyang-lissajous-Kabdelta}）恰有 \(g\) 个连通分支；每一分支同胚于 \(\TT\)，并且相位参数 \(\delta\) 仅在这些分支之间施加平移（不改变分支数）。
\end{proposition}
\begin{proof}
考虑连续群同态
\[
\phi:\TT^2\to\TT,\qquad \phi(u,v)=u^{b_0}v^{-a_0}.
\]
由于 \(\gcd(a_0,b_0)=1\)，\(\phi\) 为满射，且 \(\ker\phi\) 为一维子环面（连通）。
另一方面，约束 \(u^b=\xi v^a\) 等价于 \(\phi(u,v)^g=\xi\)，从而
\[
K_{a,b,\delta}=\bigcup_{\omega^g=\xi}\phi^{-1}(\omega).
\]
方程 \(\omega^g=\xi\) 在 \(\TT\) 上有且仅有 \(g\) 个解；而每个纤维 \(\phi^{-1}(\omega)\) 为 \(\ker\phi\) 的陪集，因而连通且同胚于 \(\TT\)。
不同 \(\omega\) 给出不交的纤维并构成连通分支分解，故分支数为 \(g\)。
参数 \(\delta\) 仅通过 \(\xi=e^{\mathrm{i}2b\delta}\) 改变上述陪集选择而不改变其个数。证毕。
\end{proof}

\begin{theorem}[迹坐标门下的双奇异环消元与有理性]\label{thm:leyang-lissajous-trace-resultant-rational}
令 \(a,b\in\ZZ_{>0}\) 与 \(\delta\in\RR\)，并考虑 Lissajous 参数曲线
\[
x(t)=\cos(at+\delta),\qquad y(t)=\cos(bt).
\]
定义逐坐标 Lee--Yang 变量
\begin{equation}\label{eq:leyang-lissajous-trace-variables}
T_x(t):=-\frac{1}{2(1+x(t))},\qquad T_y(t):=-\frac{1}{2(1+y(t))}.
\end{equation}
则由定理 \ref{thm:leyang-trace-mobius-invertibility} 有
\[
T_x(t)=J\!\left(e^{\mathrm{i}(at+\delta)}\right),\qquad
T_y(t)=J\!\left(e^{\mathrm{i}bt}\right),
\qquad (T_x(t),T_y(t))\in(-\infty,-\tfrac14]^2.
\]
令 \(z:=e^{\mathrm{i}t}\in\CC^{\times}\)，并定义关于 \(z\) 的多项式
\begin{equation}\label{eq:leyang-lissajous-trace-polynomials}
\begin{aligned}
P_x(z;X)&:=X\,e^{2\mathrm{i}\delta}z^{2a}+(2X+1)e^{\mathrm{i}\delta}z^{a}+X,\\
P_y(z;Y)&:=Y\,z^{2b}+(2Y+1)z^{b}+Y.
\end{aligned}
\end{equation}
则存在消元多项式
\[
\Phi_{a,b,\delta}(X,Y):=\Res_z\bigl(P_x(z;X),P_y(z;Y)\bigr)\in \ZZ[e^{\mathrm{i}\delta},e^{-\mathrm{i}\delta}][X,Y],
\]
使得对任意 \(t\in\RR\) 恒有
\[
\Phi_{a,b,\delta}\bigl(T_x(t),T_y(t)\bigr)=0.
\]
并且像集 \(\{(T_x(t),T_y(t)):\ t\in\RR\}\) 的 Zariski 闭包是一条不可约有理曲线，其归一化同构于 \(\PP^1\)。
\end{theorem}
\begin{proof}
由 \(X=J(v)\iff X(1+v)^2+v=0\) 等价于
\[
Xv^2+(2X+1)v+X=0.
\]
取 \(v_x=e^{\mathrm{i}(at+\delta)}=e^{\mathrm{i}\delta}z^{a}\) 与 \(v_y=e^{\mathrm{i}bt}=z^{b}\)，分别代入上式得到
\(
P_x(z;T_x(t))=0
\)
与
\(
P_y(z;T_y(t))=0
\)，其中 \(P_x,P_y\) 如 \eqref{eq:leyang-lissajous-trace-polynomials}。
因此两多项式在 \(z=e^{\mathrm{i}t}\) 处有公共根，故其结果式在 \((X,Y)=(T_x(t),T_y(t))\) 处为零，即
\(
\Phi_{a,b,\delta}(T_x(t),T_y(t))=0
\)。

其次，由 \(T_x(t)=J(e^{\mathrm{i}\delta}z^a)\)、\(T_y(t)=J(z^b)\) 可知映射
\[
\PP^1\dashrightarrow \AA^2,\qquad z\longmapsto \bigl(J(e^{\mathrm{i}\delta}z^a),\,J(z^b)\bigr)
\]
为有理映射，且非常值。
其像的 Zariski 闭包因此为不可约代数曲线；又因其被 \(\PP^1\) 有理参数化，归一化为 \(\PP^1\)，即为有理曲线。证毕。
\end{proof}

\input{thm__lee-yang-lissajous-bidegree-singularity-budget}

\begin{theorem}[特相位下的 Chebyshev 消元闭式]\label{thm:leyang-lissajous-leyang-chebyshev-elimination-special-phases}
令 \(a,b\in\NN\) 并考虑 Lissajous 族
\[
x(t)=\cos(at+\delta),\qquad y(t)=\cos(bt),
\qquad t\in\RR.
\]
定义逐坐标 Lee--Yang 逆平方变量
\[
\lambda_x(t):=-\frac{1}{4x(t)^2},\qquad \lambda_y(t):=-\frac{1}{4y(t)^2},
\]
并在 \(x(t)=0\) 或 \(y(t)=0\) 处按扩展值 \(-\infty\) 理解。
令 \(T_m\) 为第一类 Chebyshev 多项式（\(T_m(\cos\theta)=\cos(m\theta)\)）。
则：
\begin{enumerate}
  \item 当 \(\delta=0\) 时，存在单一参数 \(s=\cos(2t)\in[-1,1]\) 使
  \[
  \boxed{\;
  \lambda_x(s)=-\frac{1}{2\bigl(1+T_a(s)\bigr)},\qquad
  \lambda_y(s)=-\frac{1}{2\bigl(1+T_b(s)\bigr)}\; }.
  \]
  因而在开集 \(\{\lambda_x\lambda_y\neq 0\}\) 上满足无根号的消元关系
  \[
  \boxed{\;
  T_b\!\Bigl(-\frac{1+2\lambda_x}{2\lambda_x}\Bigr)
  =
  T_a\!\Bigl(-\frac{1+2\lambda_y}{2\lambda_y}\Bigr)\; }.
  \]
  \item 当 \(\delta=\frac{\pi}{2}\) 时（\(x(t)=-\sin(at)\)），仍存在 \(s=\cos(2t)\in[-1,1]\) 使
  \[
  \boxed{\;
  \lambda_x(s)=-\frac{1}{2\bigl(1-T_a(s)\bigr)},\qquad
  \lambda_y(s)=-\frac{1}{2\bigl(1+T_b(s)\bigr)}\; }.
  \]
  并且相应消元关系可取为
  \[
  \boxed{\;
  T_b\!\Bigl(\frac{1+2\lambda_x}{2\lambda_x}\Bigr)
  =
  T_a\!\Bigl(-\frac{1+2\lambda_y}{2\lambda_y}\Bigr)\; }.
  \]
\end{enumerate}
\end{theorem}
\begin{proof}
当 \(\delta=0\) 时，利用恒等式 \(\cos^2\alpha=\frac12(1+\cos(2\alpha))\) 得
\[
\lambda_x(t)=-\frac{1}{4\cos^2(at)}=-\frac{1}{2(1+\cos(2at))},
\qquad
\lambda_y(t)=-\frac{1}{2(1+\cos(2bt))}.
\]
令 \(s=\cos(2t)\)，则 \(\cos(2at)=T_a(s)\)、\(\cos(2bt)=T_b(s)\)，从而得到参数化闭式。
由上式反解得
\[
\cos(2at)=-\frac{1+2\lambda_x}{2\lambda_x},\qquad
\cos(2bt)=-\frac{1+2\lambda_y}{2\lambda_y},
\]
并用恒等式 \(T_b(\cos(2at))=\cos(2abt)=T_a(\cos(2bt))\) 即得消元关系。
\par
当 \(\delta=\frac{\pi}{2}\) 时有 \(x(t)=-\sin(at)\)，并由 \(\sin^2\alpha=\frac12(1-\cos(2\alpha))\) 得
\[
\lambda_x(t)=-\frac{1}{4\sin^2(at)}=-\frac{1}{2(1-\cos(2at))},
\]
其余推导同上；此时 \(\cos(2at)=\frac{1+2\lambda_x}{2\lambda_x}\) 给出第二条消元式。证毕。
\end{proof}

\paragraph{回文谱与零虚泄漏}
幺正切片上的“驻波分解/正弦模态消失”在附录 \ref{cor:trace-palindrome-completed-real} 已以回文迹振幅给出。
下述引理给出一个与该结论同型、但允许复系数的纯傅里叶判别，并将“回文谱 \(\Longleftrightarrow\) 纯余弦模态”提升为复共轭对称的必要充分条件。

\begin{lemma}[回文谱与纯余弦振幅]\label{lem:leyang-palindrome-spectrum-real-amplitude}
令 \(P(z)=\sum_{k=0}^{n}a_k z^k\) 为次数 \(n\) 的多项式（系数允许为复数），并定义
\[
A(\theta):=e^{-\mathrm{i}n\theta/2}P(e^{\mathrm{i}\theta})\qquad(\theta\in\RR).
\]
则下列条件等价：
\begin{enumerate}
  \item 对一切 \(\theta\)，有 \(A(\theta)\in\RR\)；
  \item 系数满足共轭回文条件 \(a_k=\overline{a_{n-k}}\)（\(0\le k\le n\)）。
\end{enumerate}
特别地，当 \(a_k\in\RR\) 时，该共轭回文条件化为 \(a_k=a_{n-k}\)，并推出 \(A(\theta)\) 具有严格纯余弦展开且无任何正弦模态。
\end{lemma}
\begin{proof}
展开
\[
A(\theta)=\sum_{k=0}^{n}a_k\,e^{\mathrm{i}(k-n/2)\theta},
\qquad
\overline{A(\theta)}=\sum_{k=0}^{n}\overline{a_k}\,e^{-\mathrm{i}(k-n/2)\theta}
=\sum_{k=0}^{n}\overline{a_{n-k}}\,e^{\mathrm{i}(k-n/2)\theta}.
\]
因此 \(A(\theta)=\overline{A(\theta)}\) 对所有 \(\theta\) 成立当且仅当相应 Fourier 系数逐项相等，即 \(a_k=\overline{a_{n-k}}\)。
当 \(a_k\in\RR\) 时，按 \(k\leftrightarrow n-k\) 配对立即得到纯余弦展开。证毕。
\end{proof}

\paragraph{Lee--Yang 逆平方门的二面体刻画}
坐标门 \(J(z^2)=-1/(z+z^{-1})^2\) 同时满足 \(z\mapsto z^{-1}\) 与 \(z\mapsto -z\) 的二面体不变性。
在排除非本征极点的归一化口径下，该门在二面体不变类中由其极点轨道与无穷远归一化唯一确定。

\begin{proposition}[二面体归一化下的唯一性]\label{prop:leyang-inverse-square-dihedral-uniqueness}
设 \(R(z)\) 为非常值有理函数，满足
\[
R(z)=R(z^{-1})=R(-z),
\]
并满足：
\begin{enumerate}
  \item \(R\) 在 \(\CC^{\times}\) 上仅在 \(z=\pm \mathrm{i}\) 处具有极点，且两处皆为二阶极点；
  \item \(R(1)=R(-1)=-\tfrac14\)；
  \item \(R(\infty)=0\)。
\end{enumerate}
则必有
\[
\boxed{\ R(z)=J(z^2)=-\frac{1}{(z+z^{-1})^2}\ }.
\]
\end{proposition}
\begin{proof}
由 \(R(z)=R(z^{-1})\) 与 \(R(z)=R(-z)\) 可知 \(R\) 仅依赖于 \((z+z^{-1})^2\)，从而存在有理函数 \(F\) 使
\[
R(z)=F\bigl((z+z^{-1})^2\bigr).
\]
记 \(x(z):=(z+z^{-1})^2\)。
当 \(z=\pm\mathrm{i}\) 时 \(x(z)=0\)，且 \(z\mapsto x(z)\) 在 \(\pm\mathrm{i}\) 处具有二阶零点；由 (a) 知 \(F\) 在 \(x=0\) 处为一阶极点，并且在 \(\widehat{\CC}\setminus\{0\}\) 上无其他极点。
由第三条得 \(F(\infty)=0\)，故 \(F(x)=c/x\)。
再由第二条与 \(x(1)=x(-1)=4\) 得 \(c/4=-1/4\)，即 \(c=-1\)。证毕。
\end{proof}

\paragraph{自互反多项式的 Lee--Yang 实射线化}
迹坐标 \(s=z+z^{-1}\) 把单位圆谱条件 \(|z|=1\) 压缩为区间 \(s\in[-2,2]\) 上的一元实条件；
Lee--Yang 逆平方门 \(t=J(z^2)=-1/s^2\) 进一步把该区间压缩为负实射线 \((-\infty,-\tfrac14]\)。
下述定理给出自互反多项式在该坐标门下的严格等价判别。

\begin{theorem}[自互反多项式的 Lee--Yang 实射线化判据]\label{thm:leyang-selfreciprocal-ray-criterion}
令 \(P(z)\in\RR[z]\) 为偶次数 \(2m\) 的自互反多项式：
\[
z^{2m}P(1/z)=P(z).
\]
则存在唯一 \(Q(s)\in\RR[s]\) 使
\[
P(z)=z^{m}Q(s),\qquad s:=z+z^{-1}.
\]
定义 Lee--Yang 变换
\[
\mathcal L_P(t):=t^{m}\,Q(\sqrt{-1/t})\,Q(-\sqrt{-1/t})\in\RR[t],
\]
其中乘积保证 \(\mathcal L_P\) 不依赖平方根分支且确为多项式。
则以下命题等价：
\begin{enumerate}
  \item \(P\) 的全部零点均位于单位圆 \(|z|=1\)（按重数计）；
  \item \(\mathcal L_P\) 的全部零点均位于 Lee--Yang 奇异环 \((-\infty,-\tfrac14]\)（按重数计，射线端点按紧化理解）。
\end{enumerate}
在该对应下，若 \(z\) 为 \(P\) 的单位圆零点且 \(s=z+z^{-1}\neq 0\)，则对应的 Lee--Yang 根为
\[
t=J(z^2)=-\frac{1}{s^2}\in(-\infty,-\tfrac14],
\]
与命题 \ref{prop:jouwkowsky-inverse-square} 的逆平方坐标门一致。
\end{theorem}
\begin{proof}
对合 \(\iota:z\mapsto z^{-1}\) 在 Laurent 环 \(\RR[z,z^{-1}]\) 上的不变量子环为 \(\RR[z+z^{-1}]\)，从而自互反条件等价于 \(z^{-m}P(z)\in \RR[z,z^{-1}]^{\iota}\)，即存在唯一 \(Q(s)\) 使 \(P(z)=z^{m}Q(s)\)。
令 \(s=z+z^{-1}\)，并由命题 \ref{prop:jouwkowsky-inverse-square} 得
\[
t:=J(z^2)=-\frac{1}{(z+z^{-1})^2}=-\frac{1}{s^2}.
\]
若 \(|z|=1\)，则 \(s\in[-2,2]\) 且 \(t\in(-\infty,-\tfrac14]\)。
当 \(P(z)=0\) 且 \(|z|=1\) 时，必有 \(Q(s)=0\)。
取 \(t=-1/s^2\)，则 \(\sqrt{-1/t}=\pm s\)，从而 \(\mathcal L_P(t)=0\)。
反之，若 \(t_0\in(-\infty,-\tfrac14]\) 且 \(\mathcal L_P(t_0)=0\)，令 \(s_0:=\sqrt{-1/t_0}\in[0,2]\)，则 \(Q(s_0)=0\) 或 \(Q(-s_0)=0\)。
取满足 \(z+z^{-1}=s_0\)（或 \(-s_0\)）的单位圆解 \(z\)（二次方程 \(z^2-s_0z+1=0\) 在 \(s_0\in[-2,2]\) 时有单位圆根），即得 \(|z|=1\) 且 \(P(z)=z^{m}Q(z+z^{-1})=0\)。
因此两条条件等价。证毕。
\end{proof}

\begin{corollary}[共振窗相位锁定的 Lee--Yang 锚点]\label{cor:leyang-anchor-minus4-qsqrtm15}
设
\[
z_\ast:=\frac{-1-\mathrm{i}\sqrt{15}}{4}\in\TT,
\qquad
2z_\ast^2+z_\ast+2=0.
\]
则 \(z_\ast+z_\ast^{-1}=-\tfrac12\)，从而
\[
J(z_\ast^2)=-\frac{1}{(z_\ast+z_\ast^{-1})^2}=-4.
\]
并且 \(\QQ(z_\ast)=\QQ(\sqrt{-15})\)。
该锚点与正文中共振窗影子谱包的 CM 相位锁定机制（命题 \ref{prop:pom-cm-rigidity-qsqrtm15}）一致。
\end{corollary}
\begin{proof}
由 \(2z_\ast^2+z_\ast+2=0\) 两边除以 \(z_\ast\neq 0\) 得
\(
2(z_\ast+z_\ast^{-1})+1=0
\)，故 \(z_\ast+z_\ast^{-1}=-\tfrac12\)。
代入 \(J(z^2)=-1/(z+z^{-1})^2\) 即得 \(J(z_\ast^2)=-4\)。
该二次方程的判别式为 \(-15\)，故 \(\QQ(z_\ast)=\QQ(\sqrt{-15})\)。证毕。
\end{proof}

\paragraph{单位圆测度在 Lee--Yang 奇异环上的推前密度}
命题 \ref{prop:singular-ring-limit-measure-density} 已给出由 \(t=-1/(4\cos^2\varphi)\) 推前得到的密度公式。
下述定理将该结论改写为对 \(J\) 的直接推前变换公式，便于作为积分恒等式在后续计算中调用。

\begin{theorem}[Haar 推前的严格变换公式]\label{thm:leyang-pushforward-density-formula}
令
\[
t(\theta):=J(e^{\mathrm{i}\theta})=-\frac{1}{\bigl(e^{\mathrm{i}\theta/2}+e^{-\mathrm{i}\theta/2}\bigr)^2}
=-\frac{1}{(2\cos(\theta/2))^2}\in(-\infty,-\tfrac14],
\qquad \theta\in[0,2\pi].
\]
则对任意可积函数 \(f:(-\infty,-\tfrac14]\to\CC\) 成立
\begin{equation}\label{eq:leyang-pushforward-density-integral}
\boxed{\;
\frac{1}{2\pi}\int_{0}^{2\pi} f(t(\theta))\,\dd\theta
=
\frac{1}{\pi}\int_{-\infty}^{-1/4} f(t)\,\frac{\dd t}{(-t)\sqrt{-4t-1}}\; }.
\end{equation}
\end{theorem}
\begin{proof}
令 \(\phi=\theta/2\)，则 \(\dd\theta=2\dd\phi\) 且
\[
\frac{1}{2\pi}\int_{0}^{2\pi} f(t(\theta))\,\dd\theta
=\frac{1}{\pi}\int_{0}^{\pi} f\!\left(-\frac{1}{4\cos^2\phi}\right)\dd\phi
=\frac{2}{\pi}\int_{0}^{\pi/2} f\!\left(-\frac{1}{4\cos^2\phi}\right)\dd\phi.
\]
在 \(\phi\in(0,\pi/2)\) 上作代换 \(t=-\frac{1}{4\cos^2\phi}\)。
直接计算得
\(
\frac{\dd t}{\dd\phi}=2t\sqrt{-4t-1}
\)，故
\(
\dd\phi=\frac{\dd t}{2t\sqrt{-4t-1}}
\)。
并且 \(\phi:0\mapsto t:-\tfrac14\)、\(\phi:\pi/2\mapsto t:-\infty\)。
代入并调整积分方向即得 \eqref{eq:leyang-pushforward-density-integral}。证毕。
\end{proof}

\begin{theorem}[虚轴对偶门的 Haar 推前变换公式]\label{thm:leyang-orthogonal-dual-pushforward-density-formula}
沿用命题 \ref{prop:leyang-orthogonal-dual-imaginary-projection} 的 \(J_{-}\)。
令
\[
\lambda(\theta):=J_{-}(e^{\mathrm{i}\theta})
=-\frac{1}{\bigl(e^{\mathrm{i}\theta/2}-e^{-\mathrm{i}\theta/2}\bigr)^2}
=\frac{1}{(2\sin(\theta/2))^2}\in\left[\frac14,\infty\right),
\qquad \theta\in[0,2\pi].
\]
则对任意可积函数 \(f:[\tfrac14,\infty)\to\CC\) 成立
\begin{equation}\label{eq:leyang-orthogonal-dual-pushforward-density-integral}
\boxed{\;
\frac{1}{2\pi}\int_{0}^{2\pi} f(\lambda(\theta))\,\dd\theta
=
\frac{1}{\pi}\int_{1/4}^{\infty} f(\lambda)\,\frac{\dd\lambda}{\lambda\sqrt{4\lambda-1}}\; }.
\end{equation}
\end{theorem}
\begin{proof}
令 \(\phi=\theta/2\)，则 \(\dd\theta=2\dd\phi\) 且
\[
\frac{1}{2\pi}\int_{0}^{2\pi} f(\lambda(\theta))\,\dd\theta
=\frac{1}{\pi}\int_{0}^{\pi} f\!\left(\frac{1}{4\sin^2\phi}\right)\dd\phi
=\frac{2}{\pi}\int_{0}^{\pi/2} f\!\left(\frac{1}{4\sin^2\phi}\right)\dd\phi.
\]
在 \(\phi\in(0,\pi/2)\) 上作代换 \(\lambda=\frac{1}{4\sin^2\phi}\)。
直接计算得
\(
\frac{\dd\lambda}{\dd\phi}=-2\lambda\sqrt{4\lambda-1}
\)，故
\(
\dd\phi=-\frac{\dd\lambda}{2\lambda\sqrt{4\lambda-1}}
\)。
并且 \(\phi:0\mapsto \lambda:\infty\)、\(\phi:\pi/2\mapsto \lambda:\tfrac14\)。
代入并调整积分方向即得 \eqref{eq:leyang-orthogonal-dual-pushforward-density-integral}。证毕。
\end{proof}

\begin{theorem}[Lee--Yang 反平方观测的 Haar 普适律]\label{thm:leyang-inverse-square-haar-universality}
设 \(G\) 为连通紧阿贝尔群，\(m_G\) 为其 Haar 概率测度，\(\chi:G\to\TT\) 为非平凡连续特征。
令
\[
Y(x):=J\!\bigl(\chi(x)^2\bigr)=\tau(\chi(x))\in\RR,
\]
其中 \(\tau(z)=J(z^2)\) 如 \eqref{eq:leyang-tau-def}。
则 \(Y(G)=(-\infty,-\tfrac14]\)。
并且推前测度 \(\nu:=Y_{\ast}m_G\) 在 \((-\infty,-\tfrac14)\) 上对 Lebesgue 测度绝对连续，且其密度为
\begin{equation}\label{eq:leyang-haar-universal-density}
\boxed{\;
\frac{\dd\nu}{\dd y}(y)=\frac{1}{\pi\,(-y)\,\sqrt{-4y-1}}\qquad(y\in(-\infty,-\tfrac14))\; }.
\end{equation}
\end{theorem}
\begin{proof}
由于 \(G\) 连通，\(\chi(G)\subset\TT\) 为连通闭子群；\(\chi\) 非平凡蕴含 \(\chi(G)\neq\{1\}\)，故 \(\chi(G)=\TT\)。
由 Haar 测度的唯一性，\(\chi_{\ast}m_G\) 必为 \(\TT\) 上的 Haar 概率测度 \(\lambda_{\TT}\)。
因此 \(Y(G)=\tau(\chi(G))=\tau(\TT)=(-\infty,-\tfrac14]\)（见命题 \ref{prop:jouwkowsky-inverse-square}）。
\par
对任意可积函数 \(f:(-\infty,-\tfrac14]\to\CC\)，有
\[
\int_G f(Y(x))\,\dd m_G(x)
=\int_{\TT} f(\tau(z))\,\dd\lambda_{\TT}(z).
\]
对右端应用定理 \ref{thm:leyang-pushforward-density-formula} 得
\[
\int_{\TT} f(\tau(z))\,\dd\lambda_{\TT}(z)
=\frac{1}{\pi}\int_{-\infty}^{-1/4} f(y)\,\frac{\dd y}{(-y)\sqrt{-4y-1}},
\]
从而识别出 \(\nu\) 的密度为 \eqref{eq:leyang-haar-universal-density}。证毕。
\end{proof}

\begin{corollary}[余弦相位族的“奇异环能量变量”边缘律的普适性]\label{cor:leyang-cosine-energy-marginal-universality}
令 \(\Theta\sim\mathrm{Unif}[0,2\pi)\)。
对任意整数 \(m\ge 1\) 与相位 \(\delta\in\RR\)，定义
\[
X_{m,\delta}:=\cos(m\Theta+\delta),
\qquad
T_{m,\delta}:=-\frac{1}{4X_{m,\delta}^2}\in(-\infty,-\tfrac14]\cup\{-\infty\},
\]
其中在 \(X_{m,\delta}=0\) 处按射线端点紧化取 \(T_{m,\delta}=-\infty\)（该事件概率为 \(0\)）。
则 \(T_{m,\delta}\) 的分布与 \(m,\delta\) 无关，且其在 \((-\infty,-\tfrac14)\) 上的密度恰为 \eqref{eq:leyang-haar-universal-density}。
\par
特别地，任意以 \(\cos(\cdot)\) 为径向幅度的闭轨族（例如 \eqref{eq:leyang-polar-cos-family} 的极坐标玫瑰族、以及 \eqref{eq:leyang-lissajous-family} 的 Lissajous 两坐标）在等权采样其基本相位后，
逐坐标 Lee--Yang 变量 \(-1/(4x^2)\)、\(-1/(4y^2)\) 均服从同一边缘律；所有形态差异仅由其联合分布（相关结构）承载。
\end{corollary}
\begin{proof}
因为 \(m\Theta+\delta\) 在模 \(2\pi\) 意义下仍为均匀分布，故 \(X_{m,\delta}\stackrel{d}{=}\cos\Theta\)，从而 \(T_{m,\delta}\stackrel{d}{=}-1/(4\cos^2\Theta)\)。
令 \(G=\TT\) 并取特征 \(\chi(z)=z^m\)，则 \(T_{m,\delta}=J\bigl((e^{\mathrm{i}\delta}\chi(e^{\mathrm{i}\Theta}))^2\bigr)\)。
由定理 \ref{thm:leyang-inverse-square-haar-universality}（或等价地由定理 \ref{thm:leyang-pushforward-density-formula}）即得其密度与 \(m,\delta\) 无关且等于 \eqref{eq:leyang-haar-universal-density}。证毕。
\end{proof}

\begin{corollary}[Haar 推前密度的干涉强度律]\label{cor:leyang-interference-intensity-density}
令 \(\tau\) 如 \eqref{eq:leyang-tau-def}，并记 \(\lambda_{\TT}\) 为单位圆 \(\TT\) 上的归一化 Haar 测度。
设 \(\nu:=\tau_{\ast}\lambda_{\TT}\) 为推前测度，则 \(\nu\) 在 \((-\infty,-\tfrac14]\) 上绝对连续，且其密度为
\begin{equation}\label{eq:leyang-interference-intensity-density}
\boxed{\;
\frac{\dd\nu}{\dd t}(t)=\frac{1}{\pi\,|t|\,\sqrt{4|t|-1}}\qquad(t\le -\tfrac14)\; }.
\end{equation}
\end{corollary}
\begin{proof}
对 \(z=e^{\mathrm{i}\theta}\) 有 \(\tau(z)=J(e^{\mathrm{i}2\theta})=t(2\theta)\)，其中 \(t(\cdot)\) 由 \eqref{eq:leyang-pushforward-density-integral} 给出。
由 \(t\) 的 \(2\pi\)-周期性与换元 \(\phi=2\theta\) 得
\[
\int_{\TT} f(\tau(z))\,\dd\lambda_{\TT}(z)
=\frac{1}{2\pi}\int_{0}^{2\pi} f(t(2\theta))\,\dd\theta
=\frac{1}{2\pi}\int_{0}^{2\pi} f(t(\phi))\,\dd\phi,
\]
对右端应用定理 \ref{thm:leyang-pushforward-density-formula} 即得推前密度
\(
\frac{1}{\pi}\frac{1}{(-t)\sqrt{-4t-1}}
\)；
因 \(t\le -1/4\) 时 \((-t)=|t|\) 且 \(\sqrt{-4t-1}=\sqrt{4|t|-1}\)，得到 \eqref{eq:leyang-interference-intensity-density}。证毕。
\end{proof}

\begin{corollary}[虚轴对偶门的干涉强度律]\label{cor:leyang-orthogonal-dual-interference-intensity-density}
令 \(\tau_{-}(z):=J_{-}(z^2)\)（\(z\in\TT\setminus\{\pm 1\}\)），并记 \(\lambda_{\TT}\) 为单位圆 \(\TT\) 上的归一化 Haar 测度。
设 \(\nu_{-}:=(\tau_{-})_{\ast}\lambda_{\TT}\) 为推前测度，则 \(\nu_{-}\) 在 \([\tfrac14,\infty)\) 上绝对连续，且其密度为
\begin{equation}\label{eq:leyang-orthogonal-dual-interference-intensity-density}
\boxed{\;
\frac{\dd\nu_{-}}{\dd t}(t)=\frac{1}{\pi\,t\,\sqrt{4t-1}}\qquad\left(t\ge \frac14\right)\; }.
\end{equation}
\end{corollary}
\begin{proof}
对 \(z=e^{\mathrm{i}\theta}\) 有 \(\tau_{-}(z)=J_{-}(e^{\mathrm{i}2\theta})=\lambda(2\theta)\)，其中 \(\lambda(\cdot)\) 由 \eqref{eq:leyang-orthogonal-dual-pushforward-density-integral} 给出。
由 \(\lambda\) 的 \(2\pi\)-周期性与换元 \(\phi=2\theta\) 得
\[
\int_{\TT} f(\tau_{-}(z))\,\dd\lambda_{\TT}(z)
=\frac{1}{2\pi}\int_{0}^{2\pi} f(\lambda(2\theta))\,\dd\theta
=\frac{1}{2\pi}\int_{0}^{2\pi} f(\lambda(\phi))\,\dd\phi.
\]
对右端应用定理 \ref{thm:leyang-orthogonal-dual-pushforward-density-formula} 即得推前密度 \eqref{eq:leyang-orthogonal-dual-interference-intensity-density}。证毕。
\end{proof}

\paragraph{椭圆层的正实穿越}
命题 \ref{prop:jouwkowsky-inverse-square} 指出当 \(|z|=r\neq 1\) 时，\(W(z)\) 的像为共焦椭圆。
该椭圆经逆平方门 \(J(z^2)=-1/W(z)^2\) 的像不再被束缚在负实射线上，并在同一半径层上发生正实穿越。

\begin{proposition}[离开单位圆的正实穿越判别]\label{prop:leyang-ellipse-layer-positive-real-crossing}
固定 \(r\neq 1\)，令 \(z=re^{\mathrm{i}\theta}\)。
则
\[
J(z^2)=-\frac{1}{(re^{\mathrm{i}\theta}+r^{-1}e^{-\mathrm{i}\theta})^2}
\]
的像同时与负实轴与正实轴相交，且两处交点为
\[
\theta=0:\quad J(z^2)=-\frac{1}{(r+r^{-1})^2}<0,\qquad
\theta=\frac{\pi}{2}:\quad J(z^2)=\frac{1}{(r-r^{-1})^2}>0.
\]
相反，当 \(r=1\) 时，像严格落在 Lee--Yang 奇异环 \((-\infty,-\tfrac14]\) 上。
\end{proposition}
\begin{proof}
由 \eqref{eq:jouwkowsky-inverse-square}，
\(
J(z^2)=-1/W(z)^2
\)。
当 \(\theta=0\) 时 \(W(z)=r+r^{-1}\in\RR\setminus\{0\}\)，故 \(J(z^2)=-1/(r+r^{-1})^2<0\)。
当 \(\theta=\pi/2\) 时 \(W(z)=\mathrm{i}(r-r^{-1})\)，故 \(J(z^2)=-1/(\mathrm{i}(r-r^{-1}))^2=1/(r-r^{-1})^2>0\)。
单位圆情形由命题 \ref{prop:jouwkowsky-inverse-square} 的第二条给出。证毕。
\end{proof}

\paragraph{Lissajous--奇异环泄漏的放大判别}
幺正切片上的谱穿越把临界线零点事件压缩为单位圆上同一点 \(1\) 的谱碰撞（例如推论 \ref{cor:xi-hzom-unitary-slice} 与命题 \ref{prop:xi-unitary-slice-spectral-crossing}），并在 Teichm\"uller 切片上由回文谱强制驻波型振幅与零正弦模态（结论 \ref{con:sync-kernel-weighted-teichmueller-standing-wave}；亦见引理 \ref{lem:leyang-palindrome-spectrum-real-amplitude} 的一般判别）。
另一方面，命题 \ref{prop:leyang-ellipse-layer-positive-real-crossing} 表明一旦离开单位圆，逆平方门像将不可避免地穿越正实轴，从而失去“负实射线束缚”的一维退化结构。
在该对照下，以下猜想把两通道 Lissajous 读出视为一种可放大的泄漏区分器：临界情形在去相位后退化为一维，而任何非临界偏移将迫使 Lee--Yang 坐标产生真正二维的闭轨，并可由矩放大机制给出可审计的正下界。

\begin{conjecture}[Lissajous--奇异环泄漏判别式]\label{conj:leyang-lissajous-leakage-discriminant}
在幺正切片—相位门制度下，对任一满足自对偶完成化与谱穿越口径的算术 \(\Xi\)-对象，存在自然的两通道观测（等价于在同一相位轨道上采用两组互素频率读出），使得：
\begin{enumerate}
  \item 在临界情形（对应可单位化相位与回文谱驻波约束）下，经行波相位剥离后的观测轨迹在复平面内退化为一维对象，即“虚部泄漏”严格为零；
  \item 在任意非临界偏移（对应椭圆层 \(|u|\neq 1\)）下，观测轨迹在 Lee--Yang 门像中形成真正二维的闭轨，其有向面积在适当的矩放大/多尺度粗粒化下给出可审计的严格正下界；
  \item 上述二维面积的严格正性与“非临界零点存在性”在该制度内互为充要条件，从而提供一条可放大的广义 RH 区分器。
\end{enumerate}
\end{conjecture}

\endinput

